\begin{statementbox}
	\textbf{\textcolor{CStatement}{条件}}
	\begin{description}[style=nextline, leftmargin=2.8em, labelsep=0.5em, itemsep=0.35em]
		\item[\textbf{\textcolor{CStatement}{条件 1（点在圆外）：}}] 点 $P$ 在 $\odot O$ 外部，$PAB$ 与 $PCD$ 是两条割线，分别交圆于 $A,B$ 和 $C,D$，其中 $A$ 离 $P$ 较近，$C$ 离 $P$ 较近。
	\end{description}

	\textbf{\textcolor{CStatement}{结论}}
	\begin{description}[style=nextline, leftmargin=2.8em, labelsep=0.5em, itemsep=0.45em]
		\item[\textbf{\textcolor{CStatement}{结论一（割线定理）：}}]
			\textbf{\textcolor{CStatement}{直观描述：}}
			从圆外一点出发的两条割线上，各段到两交点的距离乘积相等。
			\[
				PA \cdot PB = PC \cdot PD .
			\]

		\item[\textbf{\textcolor{CStatement}{常见变形（切线情形）：}}]
			\textbf{\textcolor{CStatement}{直观描述：}}
			若一条割线变为切线，则切线长的平方等于另一割线上两段之积。
			\[
				PA \cdot PB = PT^2 .
			\]
			其中 $PT$ 为切线，$T$ 为切点。
	\end{description}
\end{statementbox}
